\begin{abstract}

The pace of scientific research, vital for improving human life, is complex, slow, and needs specialized expertise. Meanwhile, novel, impactful research often stems from both a deep understanding of prior work, and a cross-pollination of ideas across domains and fields. To enhance the productivity of researchers, we propose ResearchAgent, which leverages the encyclopedic knowledge and linguistic reasoning capabilities of Large Language Models (LLMs) to assist them in their work. This system automatically defines novel problems, proposes methods and designs experiments, while iteratively refining them based on the feedback from collaborative LLM-powered reviewing agents. Specifically, starting with a core scientific paper, ResearchAgent is augmented not only with relevant publications by connecting information over an academic graph but also entities retrieved from a knowledge store derived from shared underlying concepts mined across numerous papers. Then, mimicking a scientific approach to improving ideas with peer discussions, we leverage multiple LLM-based ReviewingAgents that provide reviews and feedback via iterative revision processes. These reviewing agents are instantiated with human preference-aligned LLMs whose criteria for evaluation are elicited from actual human judgments via LLM prompting. We experimentally validate our ResearchAgent on scientific publications across multiple disciplines, showing its effectiveness in generating novel, clear, and valid ideas based on both human and model-based evaluation results. Our initial foray into AI-mediated scientific research has important implications for the development of future systems aimed at supporting researchers in their ideation and operationalization of novel work\footnote{Code: \href{https://github.com/JinheonBaek/ResearchAgent}{https://github.com/JinheonBaek/ResearchAgent}.}. 

\end{abstract}